\begingroup
\setlength{\emergencystretch}{2em}
\capituloTesis{Marco teórico}
\seccionTesis{Antecedentes}
\subseccionTesis{Antecedentes internacionales}

\noindent\textbf{Diniz y Almeida de Melo (2023)}

Diniz y Almeida de Melo (2023), en el artículo científico \textit{Analysis of the drainage system between the condition of urban pavements and drainage elements}, publicado en la revista \textit{ALCONPAT}, volumen 13, número 2, desarrollaron una investigación orientada a determinar la correlación entre la condición superficial del pavimento y el estado de los elementos de drenaje pluvial en el barrio de Tambaú, João Pessoa, Brasil. El problema abordado consistió en establecer si las deficiencias observadas en cunetas y sumideros mantenían una correspondencia cuantificable con el deterioro de la superficie vial, considerando que ambos componentes suelen evaluarse mediante procedimientos diferentes. Para desarrollar el estudio, seleccionaron diecinueve tramos urbanos, en los cuales determinaron el índice de condición del pavimento (PCI) mediante inspecciones superficiales, mientras que la condición de los elementos de drenaje se calificó mediante apreciación visual subjetiva; posteriormente, relacionaron ambas evaluaciones mediante análisis estadístico, con el propósito de identificar la magnitud de las asociaciones. Los resultados mostraron correlaciones débiles entre la condición del pavimento y los elementos de drenaje examinados, por lo que los autores concluyeron que las condiciones observadas no permitían establecer una correspondencia estadística fuerte entre ambos componentes. La investigación se desarrolló bajo condiciones propias de una red vial urbana, donde las características de las instalaciones pluviales, las intervenciones de mantenimiento y las condiciones particulares de cada tramo forman parte del contexto de interpretación; asimismo, los autores reconocieron que la evaluación subjetiva del drenaje constituía una característica metodológica relevante para comprender los resultados obtenidos.

Este antecedente aporta evidencia directamente relacionada con la pregunta central de la investigación en Simbal, puesto que examinó la asociación entre el drenaje superficial y la condición del pavimento, aunque empleó una calificación visual del drenaje diferente de los indicadores geométricos previstos para el corredor Pedregal–Censo–Cumbray. Su utilidad consiste en proporcionar un referente empírico para discutir las asociaciones entre la condición del drenaje longitudinal y transversal y el PCI, particularmente cuando los coeficientes obtenidos sean reducidos o no alcancen significancia estadística. La diferencia metodológica radica en que la investigación de Simbal utilizará la continuidad longitudinal libre y la proporción de abertura transversal libre como indicadores cuantitativos, además de incorporar el IRI y la correspondencia espacial entre registros, por lo que los resultados deberán compararse considerando las diferencias de medición, infraestructura y contexto territorial.

\noindent\textbf{Farida et al. (2025)}

Farida et al. (2025), en el artículo científico \textit{Integrating Hydrological Modeling and Pavement Evaluation to Analyze the Impact of Drainage Performance on Road Infrastructure}, publicado en el \textit{Journal of Geoscience, Engineering, Environment, and Technology}, volumen 10, número 4, evaluaron el sistema de drenaje y la condición del pavimento de la vía Gatot Subroto, ubicada en el distrito de Karangpawitan, Garut, Indonesia. El estudio abordó las deficiencias de una conducción pluvial abierta de sección rectangular, con aproximadamente 1 580 m de longitud y dimensiones variables, donde se habían registrado problemas de desbordamiento asociados con aportes adicionales provenientes del drenaje residencial y con secciones de capacidad insuficiente. La investigación adoptó un enfoque cuantitativo que integró análisis de frecuencia de precipitaciones, estimación de caudales de diseño, modelación hidrológica e hidráulica mediante SWMM e inspección superficial del pavimento utilizando el PCI. Para evaluar el comportamiento de los canales se consideraron periodos de retorno de dos, cinco y diez años, diferenciando las secciones que podían conducir los caudales estimados de aquellas cuya capacidad resultaba excedida. Los resultados mostraron que, para el periodo de retorno de dos años, la mayoría de las secciones presentaba condiciones adecuadas, mientras que para cinco y diez años se identificaron tramos con capacidad insuficiente; asimismo, la evaluación superficial produjo un PCI promedio de 58,66, clasificado como regular (\textit{Fair}), junto con sectores de deterioro más pronunciado. Posteriormente, los autores formularon modificaciones de las dimensiones hidráulicas y comprobaron mediante el modelo que las relaciones entre caudal máximo y capacidad a sección llena permanecían por debajo de uno en los escenarios evaluados.

Su relación con la investigación de Simbal se encuentra en la evaluación conjunta de dos componentes de la infraestructura vial: las condiciones del drenaje y el deterioro superficial del pavimento. La diferencia principal reside en la naturaleza de la variable hidráulica, puesto que Farida et al. estimaron caudales y capacidad de conducción mediante simulación, mientras que en Simbal se examinarán condiciones geométricas observables, expresadas mediante continuidad longitudinal y abertura libre de las estructuras transversales. Por consiguiente, sus resultados constituyen un referente para discutir la coexistencia de restricciones de evacuación y deterioros superficiales, sin trasladar los valores hidráulicos de Garut al corredor peruano ni interpretar una asociación estadística como demostración de causalidad. Además, el PCI promedio reportado permite establecer una comparación descriptiva, siempre que se considere la diferencia entre los sectores, las unidades de evaluación y los procedimientos utilizados.

\noindent\textbf{Winurseto et al. (2023)}

Winurseto et al. (2023), en el artículo científico \textit{Modeling of Road Performance Assessment Based on Pavement, Shoulder, and Drainage}, publicado en el \textit{Journal of Applied Engineering Science}, volumen 21, número 2, desarrollaron un modelo cuantitativo de evaluación del desempeño vial que incorporó simultáneamente la condición del pavimento, las bermas y el drenaje, con la finalidad de superar una valoración limitada exclusivamente a la superficie de rodadura. La investigación se desarrolló en Indonesia y empleó modelamiento de ecuaciones estructurales para examinar la contribución de los componentes viales al desempeño representado mediante el índice internacional de rugosidad (IRI). La operacionalización incluyó factores asociados con la condición de las bermas, tales como asentamientos, acumulación de agua, baches, diferencias de altura y pendientes, mientras que el drenaje se caracterizó mediante condiciones de la sección transversal de los canales, estado estructural y pendiente longitudinal. Los resultados atribuyeron al pavimento, las bermas y el drenaje contribuciones relativas de 58,1 \%, 20,2 \% y 21,7 \%, respectivamente, dentro del modelo desarrollado; adicionalmente, para la evaluación del drenaje se reportaron participaciones de 34,6 \% para la condición de la sección transversal, 31,1 \% para la condición estructural y 29,2 \% para la pendiente del canal. Los autores concluyeron que una evaluación del desempeño de la infraestructura vial debía considerar conjuntamente los componentes superficiales y las instalaciones de drenaje, puesto que las condiciones de estos últimos guardaban relación con el comportamiento funcional representado por la rugosidad.

Este antecedente proporciona un referente cuantitativo para relacionar las características del drenaje con el IRI, debido a que incorpora propiedades físicas de los canales dentro de un modelo de desempeño vial. Su pertinencia para Simbal se concentra en el reconocimiento de que la geometría y la condición observable del drenaje constituyen componentes diferenciables de la evaluación funcional, aunque sus indicadores no son equivalentes a la continuidad longitudinal libre ni a la abertura transversal libre propuestas en la presente investigación. La diferencia metodológica también es significativa, dado que Winurseto et al. emplearon ecuaciones estructurales para estimar contribuciones relativas, mientras que el estudio de Simbal utilizará correlaciones entre indicadores medidos por unidad estructura–sector. Por esta razón, los porcentajes del antecedente no deben interpretarse como coeficientes de correlación ni como proporciones de deterioro atribuibles al drenaje del corredor peruano.

\noindent\textbf{De Jager (2024)}

De Jager (2024), en la tesis de maestría \textit{Improvements to the Hydraulic Performance of Culverts under Inlet Control Conditions through the Optimisation of Inlet Characteristics}, desarrollada en la Universidad de Pretoria, Sudáfrica, investigó alternativas para mejorar el comportamiento hidráulico de alcantarillas existentes mediante modificaciones en sus condiciones de entrada. El problema estudiado se relacionó con las limitaciones de capacidad de las estructuras transversales y con las dificultades técnicas y económicas que implica su reemplazo, particularmente cuando una intervención requiere interrumpir la circulación vehicular o ejecutar excavaciones sobre una carretera en servicio. El autor desarrolló una investigación experimental mediante un modelo físico de alcantarilla de un solo conducto, para el cual preparó tres configuraciones de cabezales y aletas con ángulos de 90°, 45° y 30°; las dos últimas configuraciones incorporaron componentes fabricados mediante impresión tridimensional. El procedimiento incluyó la evaluación de condiciones de control de entrada, variaciones de configuración geométrica y dispositivos de ventilación, con el propósito de identificar su relación con los caudales conducidos por la estructura. Los resultados mostraron incrementos de caudal de 18 \% y 16 \% para las configuraciones de 45° y 30°, respectivamente, frente a la configuración de 90°, bajo una relación de carga hidráulica H1/D igual a 1,2 y las condiciones específicas del modelo; asimismo, el estudio identificó modificaciones del régimen de control hidráulico asociadas con la incorporación del dispositivo de ventilación. El autor concluyó que determinadas modificaciones geométricas podían mejorar el desempeño de las alcantarillas existentes, aunque su conveniencia debía examinarse considerando las condiciones de operación y los costos de intervención.

La investigación de De Jager aporta evidencia experimental sobre la importancia de las condiciones geométricas de entrada en el comportamiento de una estructura transversal, aspecto directamente relacionado con la inspección de los pases de agua del corredor Pedregal–Censo–Cumbray. Sin embargo, la diferencia metodológica debe mantenerse explícita: mientras De Jager evaluó caudales y condiciones de control hidráulico, el estudio de Simbal determinará la proporción de abertura libre observada en la entrada y la salida, sin estimar su capacidad de descarga. En consecuencia, el antecedente permite fundamentar la necesidad de identificar separadamente ambos extremos de las alcantarillas y conservar sus dimensiones originales, pero no autoriza a convertir automáticamente un porcentaje de abertura libre en un porcentaje de capacidad hidráulica. Esta distinción será particularmente relevante cuando se interpreten las asociaciones entre la condición transversal y los indicadores PCI e IRI.

\noindent\textbf{Mukwevho (2022)}

Mukwevho (2022), en la tesis de maestría \textit{The Potential Clogging and Filter Performance of Selected Geotextiles with Different Soil Types under Unidirectional Flow Conditions in Sub-soil Drainage Applications}, presentada en la Universidad de Pretoria, Sudáfrica, investigó el comportamiento de geotextiles empleados como filtros en sistemas de drenaje subsuperficial. El problema abordado consistió en determinar cómo la interacción entre las propiedades granulométricas del suelo y las características de los materiales filtrantes podía modificar la permeabilidad y favorecer procesos de obstrucción durante la circulación del agua. La investigación adoptó un enfoque experimental y utilizó ensayos de compatibilidad de filtración de larga duración para examinar cinco geotextiles en combinación con tres tipos de suelo, considerando indicadores como permeabilidad, relación de gradiente, coeficiente de uniformidad y potencial de colmatación. Los sistemas suelo–geotextil fueron sometidos a una carga hidráulica máxima de 1 420 mm durante periodos de hasta 1 008 horas, o hasta alcanzar condiciones de equilibrio. Los resultados mostraron que los suelos con proporciones elevadas de arcilla y limo favorecían la obstrucción de los poros del material filtrante, mientras que las partículas de arena y grava podían formar estructuras de retención de materiales finos. Las combinaciones de grava arenosa con los geotextiles Bidim A2 y Bidim A4 presentaron relaciones de gradiente inferiores a uno, asociadas con condiciones de filtración más abiertas, mientras que otras combinaciones registraron valores superiores a uno y mayor restricción al flujo; no obstante, ninguno de los sistemas evaluados alcanzó una condición de colmatación completa durante los ensayos.

La pertinencia de este antecedente para Simbal se encuentra en la diferenciación entre presencia de material, restricción física y capacidad de circulación del agua, aunque el estudio experimental de Mukwevho corresponde a drenaje subsuperficial y no a alcantarillas o cunetas abiertas. Sus resultados permiten comprender que la obstrucción depende de las características del material acumulado y de las condiciones del sistema, por lo que una inspección geométrica proporciona información distinta de un ensayo hidráulico. En la presente investigación, la abertura transversal libre y la continuidad longitudinal se obtendrán mediante mediciones físicas de las estructuras y recorridos observados, mientras que las características de sedimentos y residuos podrán documentarse como información descriptiva complementaria. Así, el antecedente fundamenta la necesidad de conservar la distinción entre restricción visible y desempeño hidráulico, evitando atribuir a los indicadores de Simbal propiedades que no fueron directamente medidas.
